% Gauss, Werke X/1 (1917), native118–120, printed111–113.
\SourceLeaf{118}\NoSourceHead
\begin{center}\large\textls[120]{NACHLASS UND BRIEFWECHSEL.}\end{center}
\HistoricTitle{[ÜBER DIE KREISTHEILUNGSGLEICHUNG.]}
\SourceGroup{I}
\begin{center}[ÜBER DIE PERIODEN VON $\dfrac{p-1}{3}$ UND $\dfrac{p-1}{4}$ GLIEDERN.]\end{center}
\NumberedItem{1}
\HistoricCitation{[Eintragung im \textsc{Hellwig}*), letzte Einbandseite.]}
\noindent\hspace*{\parindent}Aequationes pro summis periodorum subtriplarum ${}^{6n+1}\!\surd1$.
\[
3x=z-1
\]
\[
\begin{array}{r@{\qquad}l@{\quad}l}
7.&x^3+x^2-2x-1=z^3-21z-7&\displaystyle\sum{}^3\!\surd\frac{21\pm21\surd-3}{2}\\[8pt]
13.&x^3+x^2-4x+1=z^3-39z+65&\displaystyle\sum{}^3\!\surd\frac{39\pm39\surd-3}{2}.
\end{array}
\]
\NumberedItem{2}
\HistoricCitation{[Eintragung im \textsc{Leiste}, S. 8.]}
\noindent\hspace*{\parindent}Ad solutionem aequ: $x^p=1$ aeq: conditionales:
\begin{center}\textit{secundi gradus}\end{center}
\[
\begin{array}{r|l}
xx+x-n=0&\text{pro }p=4n+1\\
xx+x+n=0&\text{pro }p=4n-1
\end{array}
\]
\par\smallskip\noindent\rule{.25\linewidth}{.4pt}\par\smallskip
{\small\noindent\hspace*{\parindent}[*) Gemeint ist das Exemplar von J. Chr. L. \textsc{Hellwig},
\textit{Anfangsgründe der allgemeinen Mathematik u. s. w.,} Braunschweig 1777, in das \textsc{Gauss}
mehrere wissenschaftliche Bemerkungen eingetragen hat.]\par}
\clearpage
\SourceLeaf{119}\ArithmeticHead{112}{L}{ALGEBRA.\quad NACHLASS.}
\begin{center}\textit{tertii gradus}\end{center}
\begin{center}
\setlength{\arraycolsep}{5pt}\renewcommand{\arraystretch}{1.6}
\resizebox{\linewidth}{!}{$\begin{array}{r|l|l|l}
\hline\hline
p&\multicolumn{1}{c|}{\text{aequ. 1.}}&\text{aeq. 2. e prima }z=3x+1&\text{Radicum binae partes}\\
\hline\hline
7&x^3+x^2-2x-1=0&z^3-21z-7=0&{}^3\!\surd(7\pm21\surd-3):2\\
13&x^3+x^2-4x+1=0&z^3-39z+65=0&{}^3\!\surd(-65\pm39\surd-3):2\\
19&x^3+x^2-6x-7=0&z^3-57z-133=0&{}^3\!\surd(133\pm57\surd-3):2\\
31&x^3+x^2-10x-8=0&z^3-93z-4\mathpunct{.}\,31=0&\\
37&x^3+x^2-12x+7=0&z^3-111z+11\mathpunct{.}\,37=0&
\end{array}$}
\end{center}
\noindent\hspace*{\parindent}Aequationis secundae forma generalis
\par\medskip
\noindent\begin{minipage}[t]{.49\linewidth}
\[z^3-3pz-kp\]
\noindent $k$ determinatur hoc modo:\par
\noindent sit
\[p=tt+3uu\]
\noindent et erit
\[k=\pm t\pm3u.\]
\end{minipage}\hfill
\begin{minipage}[t]{.49\linewidth}
\[\begin{aligned}
\text{Rad.}&={}^3\!\surd(kp\pm l\surd-3):2\\
l&=p(t\pm u)
\end{aligned}\]
\end{minipage}
\NumberedItem{3}
\HistoricCitation{[Aus Handbuch 18, Bd, Oktober 1805, S. 147.]}
\begin{center}\textbf{Cubus.}\end{center}
\noindent\hspace*{\parindent}Es sei
\[4p=aa+27bb.\]
\noindent Man mache
\[\frac a{\surd4p}=\cos\varphi,\qquad\frac{b\surd27}{\surd4p}=\sin\varphi.\]
\noindent Dann ist der Period von $\frac13(p-1)$ Gliedern
\[=-\tfrac13+\tfrac23\cos\tfrac13\varphi\surd p.\]
\begin{center}\textbf{Biquadrat.}\end{center}
\noindent\hspace*{\parindent}Es sei
\[p=aa+4bb,\]
\[\frac a{\surd p}=\cos\varphi,\qquad\frac{2b}{\surd p}=\sin\varphi\]
\clearpage
\SourceLeaf{120}\ArithmeticHead{113}{R}{ÜBER DIE PERIODEN DER KREISTHEILUNGSGLEICHUNG.}
\noindent Der Period von $\frac14(p-1)$ Gliedern
\[\begin{aligned}
&=-\tfrac14+\tfrac14\surd p+\tfrac12\surd p\cos\tfrac12\varphi\\
&=\Pi
\end{aligned}\]
\[\Pi^0+i\Pi'-\Pi''-i\Pi'''=\surd p\mathpunct{.}\,{}^4\!\surd\frac{a+b\surd-4}{a-b\surd-4}.\]
\begin{center}\textbf{BEMERKUNG.}\end{center}
\begingroup\small
\noindent\hspace*{\parindent}Zu den Aufzeichnungen [1.] und [2.] vergleiche man die Tagebuchnotizen Nr. 39
vom 1. Okt. 1796 und Nr. 67 vom 20. Juli 1797, zu [3.] die Tagebuchnotiz Nr. 128 aus dem Jahre 1806.
Wir können hiernach die vorstehenden Aufzeichnungen als die ersten nach der Entdeckung ansehen.
Vergl. \textit{Disquisitiones Arithmeticae} art. 358 (Werke I, S. 445) und die bezüglichen Bemerkungen
im Artikel 15 des \textsc{Bachmann}schen Aufsatzes »Über \textsc{Gauss}’ zahlentheoretische Arbeiten«.
\par\noindent\hfill\textsc{Schlesinger}.\par
\endgroup
\ClosingRule
\clearpage
